
Can auxiliary training objectives teach language models to preserve user agency? Large language models aligned via Direct Preference Optimization (DPO) learn to produce responses humans prefer, but preference labels capture what users like---not necessarily what helps users maintain understanding and control. This work tests whether adding an agency-preserving auxiliary objective to DPO produces models that are both helpful and collaborative.

\subsubsection{The Agency Gap in Alignment}

Standard alignment methods optimize a single direction: making AI outputs match human preferences. DPO learns directly from comparison data indicating which response humans preferred. This framing treats the human as a passive judge rather than an active collaborator. Theoretical work suggests this may inadvertently deplete user agency---the model optimizes for immediate approval rather than long-term capability transfer \citep{mitelut2023intent}.

Three levels of this problem can be identified. At the surface level, models give confident answers without explaining reasoning. At a deeper level, preference data contains implicit signals about explanation quality, uncertainty acknowledgment, and engagement cues---but these are orthogonal to which response was preferred. At the fundamental level, the question is whether multi-objective training can capture these agency-preserving patterns and transfer them to generation behavior.

\subsubsection{Bidirectional DPO}

BiDPO (Bidirectional Direct Preference Optimization) augments standard DPO with an auxiliary agency loss:

$$\mathcal{L}_{\text{BiDPO}} = \mathcal{L}_{\text{DPO}} + \lambda \cdot \mathcal{L}_{\text{agency}}$$

The agency loss penalizes responses lacking collaborative markers---reasoning traces, uncertainty acknowledgment, and engagement invitations---quantified via a heuristic collaboration score. Unlike previous multi-objective extensions to DPO (MODPO, PAMA), this auxiliary objective targets patterns distinct from preference labels rather than additional preference dimensions.

\subsubsection{Key Findings}

The experiments reveal a partial validation with an important negative result:

\begin{enumerate}
\item \textbf{Orthogonal signal extraction works.} The collaboration score is nearly uncorrelated with preference labels (r = $-0.026)$, confirming that agency-like signals can be extracted from preference data without redundancy.
\end{enumerate}

\begin{enumerate}
\item \textbf{Multi-objective training is stable.} BiDPO training converges normally over 250 steps with zero numerical instabilities.
\end{enumerate}

\begin{enumerate}
\item \textbf{Generation-time transfer fails at proof-of-concept scale.} When evaluated on held-out prompts, BiDPO-trained models show only marginal improvement in collaboration scores over baseline DPO (+0.54\%), failing to achieve statistical significance (p = 0.247, Cohen's d = 0.016).
\end{enumerate}

The gap between training-time signal and generation-time behavior is the central finding.

\subsubsection{Contributions}

This paper makes three contributions:

\begin{enumerate}
\item Demonstration that agency-like signals exist orthogonally in preference data and can be integrated into DPO training without destabilization.
\end{enumerate}

\begin{enumerate}
\item A negative result on training-generation transfer: proof-of-concept-scale auxiliary objectives do not automatically improve downstream behavior.
\end{enumerate}

\begin{enumerate}
\item Identification of specific next steps---longer training, $\lambda$ sweeps, learned collaboration scores---that could resolve the transfer gap.
\end{enumerate}

